\clearpage
\section*{Appendix AU. Spectral population forced by terminal concentration}
\addcontentsline{toc}{section}{Appendix AU. Spectral population forced by terminal concentration}

Let \(u(t,\cdot)\) be a smooth, real-valued, periodic divergence-free
velocity on the unit three-torus at each time \(t<T\), and denote its
vector Fourier coefficients by \(\widehat u(k,t)\in\mathbb C^3\).
For any fixed \(t<T\), smoothness entails absolute convergence of
\(\sum_{k\in\mathbb Z^3}\|\widehat u(k,t)\|_{\mathbb C^3}\).
When \(E(t)=\|u(t)\|_2^2>0\), define the effective spectral population
\[
 \mathcal N_{\mathrm{eff}}(t)=
 \frac{\left(\sum_{k\in\mathbb Z^3}
       \|\widehat u(k,t)\|_{\mathbb C^3}\right)^2}
      {\sum_{k\in\mathbb Z^3}
       \|\widehat u(k,t)\|_{\mathbb C^3}^2}.
\]
This is an \(\ell^1/\ell^2\) concentration quotient, rather than a
count of modes above an arbitrary numerical cutoff.

\newtheorem{auproposition}{Proposition}
\renewcommand{\theauproposition}{AU.\arabic{auproposition}}
\begin{auproposition}[Necessary terminal spectral population]
\label{prop:AU-population}
Suppose that, for a sequence \(\tau_j\downarrow0\), the total velocity
\(u(T-\tau_j,\cdot)\) satisfies, with constants \(c,C>0\) and
\(p,\alpha>0\),
\[
 \|u(T-\tau_j)\|_2^2\le C\tau_j^p,\qquad
 \|u(T-\tau_j)\|_\infty\ge c\tau_j^{-\alpha}.
\]
Then at every time in this sequence,
\begin{equation}\label{eq:AU-population}
 \mathcal N_{\mathrm{eff}}(T-\tau_j)\ge
  \frac{c^2}{C}\,\tau_j^{-(2\alpha+p)}.
\end{equation}
If the Fourier series has finite support \(S_j\) at these times, its
cardinality must satisfy
\(
 |S_j|\ge(c^2/C)\tau_j^{-(2\alpha+p)}.
\)
\end{auproposition}

\begin{proof}
The triangle inequality, absolute Fourier convergence, and Parseval give
\[
 c\tau_j^{-\alpha}
 \le \|u(T-\tau_j)\|_\infty
 \le\sum_k\|\widehat u(k,T-\tau_j)\|,
 \qquad
 \sum_k\|\widehat u(k,T-\tau_j)\|^2
 =\|u(T-\tau_j)\|_2^2
 \le C\tau_j^p.
\]
Squaring the former and dividing by the latter yields
\eqref{eq:AU-population}. If the support is finite,
Cauchy--Schwarz also yields
\[
 \left(\sum_{k\in S_j}\|\widehat u(k)\|\right)^2
 \le |S_j|\sum_{k\in S_j}\|\widehat u(k)\|^2,
\]
so \(\mathcal N_{\mathrm{eff}}\le |S_j|\).
\end{proof}

For a terminal carrier whose squared energy is comparable to
\(\tau^{1/5}\), the same upper energy bound on the \emph{full}
corrected velocity and a retained lower speed
\(\|u(T-\tau_j)\|_\infty\ge c\tau_j^{-(1/2+h)}\)
force
\begin{equation}\label{eq:AU-exponent}
 \mathcal N_{\mathrm{eff}}(T-\tau_j)
 \gtrsim\tau_j^{-(6/5+2h)}.
\end{equation}
The illustrative \(h=1/10\) class requires exponent \(7/5\).
Neither assumption follows from the 20-mode Fourier calculation
in Appendix~AT; both refer to the full, unforced correction.
Equation~\eqref{eq:AU-exponent} is therefore a necessary
spectral-concentration test, not a construction.

\begin{auproposition}[Necessary frequency radius under a global band limit]
\label{prop:AU-band}
Under the hypotheses of Proposition~\ref{prop:AU-population},
suppose additionally that the Fourier support of the \emph{full}
velocity at the terminal times lies in
\(
 \{k\in\mathbb Z^3:\|k\|_\infty\le K_0\tau_j^{-\beta}\},
\)
for constants \(K_0,\beta>0\). Then necessarily
\[
 \beta\ge\frac{2\alpha+p}{3}.
\]
For \(p=1/5\) and \(\alpha=3/5\), this is
\(\beta\ge7/15\). A global radius exponent
\(\beta=9/20<7/15\) is excluded.
\end{auproposition}

\begin{proof}
The Fourier support has at most
\((2\lfloor K_0\tau_j^{-\beta}\rfloor+1)^3
=O(\tau_j^{-3\beta})\) elements.
By Cauchy--Schwarz this bound also controls
\(\mathcal N_{\mathrm{eff}}\).
Together with \eqref{eq:AU-population}, it implies
\(2\alpha+p\le 3\beta\) as \(\tau_j\downarrow0\).
\end{proof}

In particular, a frequency-migrating exterior field with ceiling
\(O(\tau^{-9/20})\) cannot by itself serve as the \emph{entire}
velocity subject to both the stated energy and speed scalings.
The shrinking carrier has higher spatial frequencies, so this
conditional band-limit obstruction does \emph{not} rule out a
combined carrier--halo construction. The bound is an elementary
Fourier/Bernstein consequence, not a dynamical transfer theorem.

The finite Fourier inequality underlying the proof is also encoded
in \path{UnforcedPatching/Attack25SpectralPopulation.lean}.
That finite-sum theorem does not by itself formalize the infinite
Fourier-series reconstruction or the specific carrier asymptotics;
those remain separate physical and analytic verification
requirements. The general inequalities in this appendix are
standard consequences of Fourier reconstruction and
Cauchy--Schwarz and are not presented as new prior-art claims.
